\section{Repository audit and a path to Lean formalization}
\label{sec:repository-audit}

\subsection{The exact first-order theory already available}

The repository was inspected at the immutable commit
\begin{center}
\texttt{22afac4c9db7a0e060c081cc11db0dc96fe59e93}.
\end{center}
The toolchain specifies Lean~4.32.0. Its Mathlib revision is
\begin{center}
\texttt{81a5d257c8e410db227a6665ed08f64fea08e997}.
\end{center}
The root Lake configuration declares the library \texttt{ConciseIPRA}
with source directory
\path{Logic/PrimitiveRecursiveArithmetic/ConciseIPRA/Lean}.
These are observations of the files, not claims that this session rebuilt
the repository.\footnote{\href{https://github.com/VladimirReshetnikov/ProveIt/tree/22afac4c9db7a0e060c081cc11db0dc96fe59e93/Logic/PrimitiveRecursiveArithmetic/ConciseIPRA}{Pinned ConciseIPRA source tree}; the same commit's root
\texttt{lean-toolchain}, \texttt{lake-manifest.json}, and
\texttt{lakefile.toml} supply the toolchain and build configuration.}

The inspected source files supply the following interfaces.
\begin{center}
\small
\begin{tabular}{@{}>{\raggedright\arraybackslash}p{0.38\linewidth}>{\raggedright\arraybackslash}p{0.56\linewidth}@{}}
\hline
Lean module & Relevant definitions and results\\
\hline
\texttt{Syntax.lean} & Scoped first-order formulas \texttt{Fm}; the predicate
\texttt{Fm.QF} for formulas with no quantifiers.\\
\texttt{Deduction.lean} & The intuitionistic natural-deduction judgment
\texttt{Prf Ax Gamma phi}, including equality, substitution, and weakening.\\
\texttt{PR.lean} & Primitive-recursive symbols \texttt{PR}, terms
\texttt{PTm}, axioms \texttt{PRAAx}, and the internal quantifier-free
induction rule \texttt{PRA.qfInd}.\\
\texttt{TeqRec.lean} & Primitive recursion derived from the scalar
iterator by a paired counter-and-value accumulator.\\
\texttt{NativeBiInterpretation.lean} & Faithful translations
\texttt{pra\_iff\_native} and \texttt{native\_iff\_pra}, both round trips,
and equiconsistency.\\
\hline
\end{tabular}
\end{center}
All references in this subsection are to the source tree at the pinned
commit; the accompanying \texttt{lean/README.md} provides its direct URL.

The symbol constructors are zero, successor, projections, composition,
and primitive recursion.  The axioms are their defining equations,
successor separation from zero, successor injectivity, and quantifier-free
induction.  The deduction system has the introduction and elimination
rules for intuitionistic conjunction, disjunction, implication, and both
quantifiers, together with equality and elimination of falsehood.  Thus the
target theory has a concrete representation; there is no need to substitute
an informal appeal to ``constructive Lean'' for an internal arithmetic
derivation.

\subsection{What can be reused, and what still has to be constructed}

The repository already contains a package at
\begin{center}
\texttt{docs/proofs/geometry/isoperimetric\_ipra/}
\end{center}
under \texttt{Logic/PrimitiveRecursiveArithmetic}.  Its README describes
a sharp planar argument based on dyadic roots of unity, finite Fourier
identities, explicitly bounded polygonal refinement, and effective limit
presentations.  It also expressly labels its Lean files an uncompiled
blueprint and identifies the internal proof as unfinished.  Its presence
therefore supplies prior mathematics and specifications, not an already
available Lean theorem proving the inequality in iPRA.  The present report
states and proves the finite lemmas it uses, and keeps the same distinction
between a human proof, an executable diagnostic, and a completed internal
formal derivation.

The bridge \texttt{pra\_iff\_native} has the logical form
\[
 \operatorname{Prf}(\mathrm{PRAAx},\Gamma,\varphi)
 \quad\longleftrightarrow\quad
 \operatorname{NPrf}(\Gamma^{\mathrm{native}},
                         \varphi^{\mathrm{native}}).
\]
It transfers an \emph{already constructed} arithmetic derivation.  It does
not turn a Mathlib theorem about real numbers, an externally proved Boolean
specification, or a passing Python computation into that derivation.
This distinction is essential for the constrained-foundation requirement.

\subsection{Finite encodings and theorem specifications}

The supplied \texttt{lean/Outline.lean} defines rational polygon vertices,
shoelace area, edge energy, rational upper bounds on edge lengths,
radical-free triangle targets, the even and odd Fourier deficits of support
functions, rational boxes and their Minkowski sums, cube vertex masks,
internal and boundary edge counts, and closed neighborhoods.  It exposes
the one-sided parallel-volume modulus needed for the higher-dimensional
consequence.  A supplied modulus is not a theorem asserting that an
arbitrary set has such a modulus.

For internalization, encode signed integers by parity and rationals by a
signed numerator and a positive denominator.  Rational order is checked
by cross multiplication.  One concrete injective pairing is
\[
 \langle a,b\rangle=(a+b)^2+a.
\]
Its inverse on its image is found by a bounded scan over pairs at most
the input code.  Use code~$0$ for the empty list and
$1+\langle a,w\rangle$ for a list with head~$a$ and tail code~$w$.
Every tail code is smaller; bounded decoding with fuel equal to the input
code is primitive recursive.  Arrays of rationals, polygons, coefficient
lists, and box lists thereby have explicit finite numerical codes.
Cube sets are natural-number bit masks, read only at vertices below $2^d$.

For each theorem, let $C$ be its fixed primitive-recursive checker, returning
zero when the stated implication holds.  Invalid inputs or false hypotheses
return zero.  The actual internal goal is
\[
 \operatorname{Prf}\bigl(\mathrm{PRAAx},\varnothing,
                         \forall z\,[C(z)=0]\bigr).
\]
The compiler must produce a particular symbol of type \texttt{PR 1},
prove its agreement with the specified checker, and supply its defining
equations for use inside the theory.  Returning a constant zero function
does not meet the specification.  The external semantics of \texttt{PR}
syntax used in the outline is a Lean definition; it is not asserted to be
a primitive-recursive universal interpreter for every primitive-recursive
program.

\subsection{The induction obligations}

Every use of \texttt{PRA.qfInd} must include a proof that its formula is
quantifier-free.  A bounded universal is replaced by a primitive-recursive
Boolean fold; a bounded existential is replaced by bounded search with
a success flag and a witness.  Prefix tables provide the finite conjunction
of earlier stages when a proof is presented as strong induction.  These
tables have explicit bounds.  The induction invariant becomes an equation
between primitive-recursive terms.  Unbounded universal closure is introduced
only after the parameterized internal derivation is complete.

For example, the cube proof at dimension~$d$ can enumerate every mask below
$2^{2^d}$ within a Boolean fold.  Compression is iterated for the stated
rank-potential bound, with identity steps after saturation.  For the box
proof, the finite cut tree must carry the box-count measure that decreases
in each child, and the volume-ratio and disjointness invariants.  Record the
actual cut tree by a bounded queue loop and traverse it bottom-up; the
invariant is the Boolean conjunction of the inequalities at its processed
nodes.  There is no induction formula quantifying over all unbounded
rational input codes.  For polygons, all refinement
counts, interval precisions, and coefficient-error bounds must be retained.
Merely observing that a numerical routine terminates does not prove the
invariant required of it.

\subsection{Completion order and verification standard}

An implementable module order is: numerical codecs; internal ordered
arithmetic; the fixed constant-enclosure algorithms actually used; polygon
identities and finite spectral estimates; triangle and support-function
identities; box dissection and matched-cut Brunn--Minkowski; cube counting
and compressions; and finally the transfer lemmas for each fixed effective
curve or region presentation.  Every fixed presentation carries its own
primitive-recursive moduli and validity proofs.  No theorem is inferred
about all real sets or all rectifiable curves from the existence of those
particular data.

The last stage is to construct the internal checker theorems and then apply
\texttt{pra\_iff\_native}.  The formalization should publish the actual
closed \texttt{Prf PRAAx} targets, successful build output, and
\texttt{\#print axioms} reports for the ambient Lean declarations.  Neither
the target definitions nor the conditional translation theorem are reported
as completed proofs here.  No Lean or Lake executable was available on the
execution path, so the included outline has not been compiled.
